\documentclass[11pt,letterpaper]{article}

\usepackage[utf8]{inputenc}
\usepackage[T1]{fontenc}
\usepackage[spanish,es-tabla]{babel}
\usepackage{geometry}
\usepackage{booktabs}
\usepackage{longtable}
\usepackage{array}
\usepackage{graphicx}
\usepackage{float}
\usepackage{xcolor}
\usepackage{hyperref}
\usepackage{amsmath}
\usepackage{amssymb}

\geometry{margin=2.25cm}
\graphicspath{
    {../output_cluster_analisis/plots_2d_ica/}
    {output_cluster_analisis/plots_2d_ica/}
    {../output_cluster_analisis/paradoxical_analysis/figures/by_subject/}
    {output_cluster_analisis/paradoxical_analysis/figures/by_subject/}
}
\hypersetup{
    colorlinks=true,
    linkcolor=blue,
    urlcolor=blue
}

\newcommand{\code}[1]{\texttt{#1}}
\newcommand{\pct}[1]{#1\%}

\title{\textbf{Informe especializado: MAT1022}\\
\large Calculo I}
\author{Proyecto Cluster Analisis}
\date{\today}

\begin{document}

\maketitle

\begin{abstract}
Este informe presenta una lectura especializada de los resultados para \code{MAT1022}. El objetivo es identificar alumnos con porcentajes bajos en \code{Porcentaje\_DMU} y \code{Porcentaje\_GA\_GB}, pero con calificaciones altas en la materia, y despues localizar los profesores asociados. En esta materia los resultados requieren mas cautela que en \code{MAT1012}: el cluster original encuentra un grupo muy pequeno y el analisis binario GMM selecciona un grupo demasiado amplio.
\end{abstract}

\section{Objetivo de la materia}

Para cada alumno de \code{MAT1022} con datos completos se construyo el vector:

\[
X_i =
\left(
\text{Porcentaje\_DMU}_i,\,
\text{Porcentaje\_GA\_GB}_i,\,
\text{CALIFICACION}_i
\right).
\]

El patron de interes es:

\begin{center}
\textbf{porcentajes bajos en examenes de referencia, junto con calificacion alta en la materia.}
\end{center}

Se aplicaron dos lecturas complementarias. La primera es el clustering original, que filtra alumnos con \code{CALIFICACION < 7.5}. La segunda es el analisis binario/paradojico, que trabaja con todos los casos completos y compara el GMM binario, el score de discrepancia y el benchmark manual \code{40/40/8}.

\section{Clustering original con filtro de calificacion}

El clustering original usa solo observaciones completas y excluye alumnos con \code{CALIFICACION < 7.5}. Luego estandariza las variables y ajusta una mezcla gaussiana. El cluster objetivo se selecciona maximizando:

\[
S_c =
z(\text{CALIFICACION})_c
- z(\text{Porcentaje\_DMU})_c
- z(\text{Porcentaje\_GA\_GB})_c.
\]

\begin{table}[H]
\centering
\caption{Resumen del clustering original para MAT1022.}
\begin{tabular}{lr}
\toprule
Indicador & Valor \\
\midrule
Registros totales de la materia & 6,946 \\
Casos completos en $R^3$ & 2,510 \\
Casos excluidos por \code{CALIFICACION < 7.5} & 134 \\
Casos usados para clustering & 2,376 \\
Metodo & GMM \\
Numero de clusters seleccionado & 11 \\
Cluster objetivo & 6 \\
Alumnos en cluster objetivo & 71 \\
Proporcion del cluster objetivo & 2.99 \\
Score del cluster objetivo & 2.617 \\
Silhouette & 0.092 \\
Calinski-Harabasz & 438.387 \\
Davies-Bouldin & 2.094 \\
\bottomrule
\end{tabular}
\end{table}

\begin{table}[H]
\centering
\caption{Centroide del cluster objetivo de MAT1022 en escala original.}
\begin{tabular}{lrr}
\toprule
Variable & Media materia filtrada & Centroide cluster objetivo \\
\midrule
\code{Porcentaje\_DMU} & 65.31 & 42.63 \\
\code{Porcentaje\_GA\_GB} & 67.29 & 24.54 \\
\code{CALIFICACION} & 9.21 & 8.72 \\
\bottomrule
\end{tabular}
\end{table}

\subsection{Interpretacion}

En \code{MAT1022}, el cluster objetivo del analisis original tiene porcentajes muy bajos, especialmente en \code{GA-GB}, pero su centroide de \code{CALIFICACION} queda por debajo de la media filtrada de la materia. Por esta razon, el propio pipeline marca una advertencia metodologica: el cluster seleccionado maximiza el score por tener porcentajes muy bajos, pero no satisface claramente la condicion de calificacion por encima de la media.

Esto no significa que el resultado no sirva. Significa que, para \code{MAT1022}, el patron encontrado es mas bien ``porcentajes muy bajos con calificacion aprobatoria/relativamente alta'', no necesariamente ``calificacion alta respecto al promedio de la materia''. La presentacion de resultados debe incluir esa salvedad.

\section{Profesores asociados al cluster objetivo}

La siguiente tabla lista los profesores con mayor proporcion de alumnos en el cluster objetivo del clustering original. La proporcion se calcula como:

\[
\text{share}_{p} =
\frac{\#\text{ alumnos del profesor }p\text{ en cluster objetivo}}
{\#\text{ alumnos clusterizados del profesor }p}.
\]

\begin{table}[H]
\centering
\caption{Top profesores MAT1022 segun cluster objetivo original.}
\begin{tabular}{lrrrrrr}
\toprule
\code{CLAVEPROFESOR} & Total & Objetivo & Share & Cal. obj. & DMU obj. & GA-GB obj. \\
\midrule
23823 & 41 & 4 & 9.76 & 8.50 & 41.67 & 18.75 \\
22114 & 41 & 4 & 9.76 & 8.90 & 46.67 & 25.00 \\
23442 & 23 & 2 & 8.70 & 8.45 & 36.67 & 29.17 \\
23837 & 26 & 2 & 7.69 & 8.25 & 53.33 & 16.67 \\
21192 & 56 & 4 & 7.14 & 8.93 & 36.67 & 23.44 \\
23453 & 32 & 2 & 6.25 & 7.90 & 56.67 & 9.38 \\
23897 & 115 & 7 & 6.09 & 9.33 & 39.05 & 23.81 \\
12516 & 18 & 1 & 5.56 & 9.00 & 60.00 & 16.67 \\
\bottomrule
\end{tabular}
\end{table}

Los shares son mucho menores que en \code{MAT1012}. Esto es esperable porque el cluster objetivo de \code{MAT1022} es pequeno: solo 71 alumnos. Por lo tanto, cualquier ranking de profesores es sensible a pocos casos y debe revisarse junto con los tamanos muestrales.

\section{Analisis binario/paradojico}

El segundo analisis trabaja con todos los alumnos completos de \code{MAT1022}, sin filtrar previamente por \code{CALIFICACION > 7.5}. El metodo principal ajusta una mezcla gaussiana de dos componentes sobre las variables estandarizadas dentro de la materia. Tambien se compara contra un score individual de discrepancia y contra el benchmark manual:

\[
\text{DMU} < 40,\quad \text{GA-GB} < 40,\quad \text{CALIFICACION} > 8.
\]

\begin{table}[H]
\centering
\caption{Tamano de grupo objetivo por metodo en MAT1022.}
\begin{tabular}{lrr}
\toprule
Metodo & Alumnos seleccionados & Proporcion \\
\midrule
GMM binario principal & 2,362 & 94.10 \\
Score de discrepancia & 2,445 & 97.41 \\
Benchmark 40/40/8 & 39 & 1.55 \\
\bottomrule
\end{tabular}
\end{table}

\begin{table}[H]
\centering
\caption{Medias del grupo principal GMM contra el resto en MAT1022.}
\begin{tabular}{lrrr}
\toprule
Variable & Grupo GMM & Resto & Diferencia \\
\midrule
\code{Porcentaje\_DMU} & 65.31 & 53.74 & 11.58 \\
\code{Porcentaje\_GA\_GB} & 68.11 & 43.19 & 24.92 \\
\code{CALIFICACION} & 9.20 & 4.96 & 4.24 \\
\bottomrule
\end{tabular}
\end{table}

\subsection{Lectura del analisis binario}

En \code{MAT1022}, el GMM binario principal selecciona 94.10\% de los casos completos. Esta es una senal de que la particion binaria no esta encontrando un grupo objetivo pequeno y sustantivamente especifico. Aunque el grupo seleccionado tiene calificaciones mucho mas altas que el resto, tambien tiene porcentajes promedio mas altos en \code{DMU} y \code{GA-GB}. Por lo tanto, no representa el patron buscado de ``porcentajes bajos y calificacion alta''.

El benchmark \code{40/40/8} selecciona 39 alumnos, equivalente al 1.55\% de los casos completos. Aunque este criterio es arbitrario y no debe ser el metodo principal, en \code{MAT1022} funciona como una referencia estricta para localizar casos extremos de porcentajes bajos y calificacion mayor a 8.

\begin{table}[H]
\centering
\caption{Solapamiento entre metodos en MAT1022.}
\begin{tabular}{lrrr}
\toprule
Comparacion & Interseccion & Jaccard & Coincidencia \\
\midrule
GMM vs. score & 2,362 & 0.966 & 0.967 \\
GMM vs. benchmark & 32 & 0.014 & 0.069 \\
Score vs. benchmark & 39 & 0.016 & 0.041 \\
\bottomrule
\end{tabular}
\end{table}

El alto solapamiento entre GMM y score se debe a que ambos metodos seleccionan casi toda la materia. El bajisimo Jaccard contra el benchmark confirma que el grupo estadistico amplio y el grupo extremo \code{40/40/8} no estan midiendo lo mismo.

\section{Profesores bajo el analisis binario}

\begin{table}[H]
\centering
\caption{Top profesores MAT1022 segun grupo principal GMM binario.}
\begin{tabular}{lrrrrr}
\toprule
\code{CLAVEPROFESOR} & Total & Grupo GMM & Share GMM & Benchmark & Share bench. \\
\midrule
23425 & 63 & 63 & 100.00 & 3 & 4.76 \\
25503 & 22 & 22 & 100.00 & 0 & 0.00 \\
25531 & 10 & 10 & 100.00 & 0 & 0.00 \\
23885 & 125 & 124 & 99.20 & 2 & 1.60 \\
24097 & 64 & 63 & 98.44 & 0 & 0.00 \\
23978 & 49 & 48 & 97.96 & 0 & 0.00 \\
23153 & 37 & 36 & 97.30 & 0 & 0.00 \\
23434 & 33 & 32 & 96.97 & 0 & 0.00 \\
\bottomrule
\end{tabular}
\end{table}

Esta tabla no debe usarse como ranking fuerte de profesores para el patron buscado, porque el grupo GMM binario incluye casi todos los alumnos. Para \code{MAT1022}, es mas informativo revisar el cluster original pequeno y el benchmark estricto como sensibilidad, siempre reportando los tamanos pequenos y la advertencia de que el centroide del cluster original no supera la media de calificacion de la materia filtrada.

\section{Figuras}

\begin{figure}[H]
\centering
\includegraphics[width=0.9\textwidth]{MAT1022_ica.png}
\caption{Proyeccion ICA del clustering original para MAT1022. Los puntos estan coloreados por cluster. La figura ayuda a ubicar visualmente el cluster objetivo, pero la interpretacion debe recordar que el centroide de calificacion del cluster objetivo no queda por encima de la media filtrada.}
\end{figure}

\begin{figure}[H]
\centering
\includegraphics[width=0.9\textwidth]{subject_MAT1022_binary_scatter.png}
\caption{Scatter de \code{Porcentaje\_DMU} contra \code{Porcentaje\_GA\_GB} para MAT1022. El analisis binario principal no separa un grupo pequeno de bajos porcentajes; mas bien deja a la gran mayoria en el grupo GMM seleccionado.}
\end{figure}

\begin{figure}[H]
\centering
\includegraphics[width=0.9\textwidth]{subject_MAT1022_pca_scatter.png}
\caption{Proyeccion PCA del analisis binario para MAT1022. La figura muestra que la particion de dos componentes esta dominada por la separacion entre un grupo mayoritario de calificaciones altas y un resto pequeno, no por el patron especifico de bajos porcentajes con calificacion alta.}
\end{figure}

\begin{figure}[H]
\centering
\includegraphics[width=0.9\textwidth]{subject_MAT1022_group_boxplots.png}
\caption{Boxplots por grupo binario en MAT1022. El grupo GMM tiene calificaciones mucho mayores que el resto, pero tambien porcentajes mayores en \code{DMU} y \code{GA-GB}. Esta figura explica por que el grupo binario no debe interpretarse como el grupo objetivo sustantivo de bajos porcentajes y calificacion alta.}
\end{figure}

El grafico 3D interactivo de esta materia esta disponible en:

\begin{center}
\parbox{0.9\textwidth}{\centering\footnotesize\path{output_cluster_analisis/paradoxical_analysis/figures/by_subject/subject_MAT1022_binary_3d.html}}
\end{center}

\section{Conclusion especifica para MAT1022}

Para \code{MAT1022}, los resultados son mas delicados que para \code{MAT1012}. El clustering original encuentra un grupo pequeno de 71 alumnos con porcentajes muy bajos, especialmente en \code{GA-GB}, pero su calificacion media queda por debajo de la media filtrada de la materia. El analisis binario GMM, por otro lado, selecciona casi toda la materia y por eso no es util como grupo objetivo estrecho.

La recomendacion para presentar \code{MAT1022} es enfatizar que el patron existe de forma mas debil y menos estable. Los profesores listados deben considerarse puntos de revision descriptiva, no conclusiones. Para esta materia conviene complementar con inspeccion manual de alumnos, sesiones, cohortes y el benchmark \code{40/40/8}.

\IfFileExists{reportes/apendice_tablas_profesores_MAT1022.tex}{\clearpage

\appendix

\section{Tablas completas de profesores para MAT1022}

Este apendice se genera automaticamente desde \code{data/datos\_procesados/professor\_appendix\_all\_years} y \code{data/datos\_procesados/professor\_appendix\_by\_period}. El desglose temporal usa exactamente las combinaciones observadas de \code{anio} y \code{CLAVESESION}.

CÁLCULO I

\subsection{Apendice A: tabla global de profesores}

\begingroup
\scriptsize
\setlength{\tabcolsep}{3pt}
\begin{longtable}{lrrrrrrrrr}
\caption{Tabla completa de profesores para MAT1022, todos los anios.}\label{tab:prof-global-mat1022-especializado}\\
\toprule
Profesor & Total & Obj. & Obj. \% & Bench. & Bench. \% & Cal. & DMU & GA-GB & Rank \\
\midrule
\endfirsthead
\toprule
Profesor & Total & Obj. & Obj. \% & Bench. & Bench. \% & Cal. & DMU & GA-GB & Rank \\
\midrule
\endhead
23425 & 63 & 63 & 100.0 & 3 & 4.8 & 9.30 & 60.32 & 64.52 & 1 \\
25503 & 22 & 22 & 100.0 & 0 & 0.0 & 9.69 & 67.27 & 67.42 & 2 \\
25531 & 10 & 10 & 100.0 & 0 & 0.0 & 9.87 & 70.67 & 71.67 & 3 \\
23885 & 125 & 124 & 99.2 & 2 & 1.6 & 9.20 & 65.44 & 69.00 & 4 \\
24097 & 64 & 63 & 98.4 & 0 & 0.0 & 9.29 & 72.71 & 63.57 & 5 \\
23978 & 49 & 48 & 98.0 & 0 & 0.0 & 9.35 & 72.65 & 66.54 & 6 \\
23153 & 37 & 36 & 97.3 & 0 & 0.0 & 9.24 & 69.91 & 66.55 & 7 \\
23434 & 33 & 32 & 97.0 & 0 & 0.0 & 8.46 & 57.78 & 65.03 & 8 \\
21262 & 29 & 28 & 96.6 & 0 & 0.0 & 8.26 & 71.26 & 68.25 & 9 \\
20792 & 56 & 54 & 96.4 & 1 & 1.8 & 8.82 & 67.26 & 74.89 & 10 \\
23149 & 28 & 27 & 96.4 & 0 & 0.0 & 9.13 & 65.00 & 72.10 & 11 \\
24085 & 133 & 128 & 96.2 & 3 & 2.3 & 9.84 & 61.80 & 68.36 & 12 \\
23852 & 158 & 152 & 96.2 & 4 & 2.5 & 9.07 & 67.26 & 66.65 & 13 \\
24827 & 26 & 25 & 96.2 & 0 & 0.0 & 9.29 & 73.59 & 72.04 & 14 \\
23442 & 24 & 23 & 95.8 & 1 & 4.2 & 9.07 & 62.22 & 61.46 & 15 \\
23897 & 119 & 114 & 95.8 & 4 & 3.4 & 9.36 & 65.15 & 66.82 & 16 \\
22446 & 112 & 107 & 95.5 & 1 & 0.9 & 9.14 & 64.17 & 65.94 & 17 \\
23452 & 110 & 105 & 95.5 & 2 & 1.8 & 9.32 & 69.70 & 65.70 & 18 \\
23823 & 44 & 42 & 95.5 & 3 & 6.8 & 8.68 & 62.42 & 58.66 & 19 \\
23148 & 64 & 61 & 95.3 & 1 & 1.6 & 8.66 & 59.69 & 65.33 & 20 \\
23411 & 121 & 115 & 95.0 & 0 & 0.0 & 9.07 & 67.55 & 69.80 & 21 \\
21853 & 40 & 38 & 95.0 & 0 & 0.0 & 8.90 & 67.00 & 69.84 & 22 \\
12516 & 19 & 18 & 94.7 & 0 & 0.0 & 8.90 & 52.28 & 60.75 & 23 \\
22473 & 56 & 53 & 94.6 & 0 & 0.0 & 9.16 & 62.02 & 63.95 & 24 \\
23453 & 34 & 32 & 94.1 & 1 & 2.9 & 8.20 & 56.86 & 70.22 & 25 \\
23836 & 166 & 155 & 93.4 & 4 & 2.4 & 8.98 & 65.78 & 66.29 & 26 \\
22114 & 44 & 41 & 93.2 & 2 & 4.5 & 8.81 & 52.88 & 59.52 & 27 \\
23837 & 28 & 26 & 92.9 & 0 & 0.0 & 8.38 & 50.00 & 54.91 & 28 \\
23839 & 55 & 51 & 92.7 & 0 & 0.0 & 8.47 & 59.03 & 59.92 & 29 \\
24086 & 68 & 62 & 91.2 & 0 & 0.0 & 8.85 & 69.31 & 72.49 & 30 \\
23096 & 22 & 20 & 90.9 & 0 & 0.0 & 8.79 & 58.18 & 65.91 & 31 \\
21192 & 60 & 54 & 90.0 & 2 & 3.3 & 9.14 & 66.11 & 61.98 & 32 \\
18269 & 36 & 32 & 88.9 & 1 & 2.8 & 8.62 & 71.11 & 75.64 & 33 \\
23152 & 61 & 54 & 88.5 & 0 & 0.0 & 8.41 & 66.78 & 68.37 & 34 \\
25843 & 23 & 20 & 87.0 & 0 & 0.0 & 8.01 & 71.30 & 63.59 & 35 \\
22521 & 30 & 26 & 86.7 & 0 & 0.0 & 8.23 & 66.89 & 75.28 & 36 \\
21779 & 97 & 84 & 86.6 & 1 & 1.0 & 8.38 & 69.62 & 70.15 & 37 \\
23835 & 73 & 63 & 86.3 & 0 & 0.0 & 8.31 & 57.72 & 64.98 & 38 \\
22471 & 39 & 33 & 84.6 & 0 & 0.0 & 7.75 & 59.66 & 64.74 & 39 \\
23874 & 25 & 21 & 84.0 & 0 & 0.0 & 7.83 & 50.40 & 63.75 & 40 \\
23162 & 6 & 6 & 100.0 & 1 & 16.7 & 8.93 & 54.44 & 55.56 &  \\
23887 & 3 & 3 & 100.0 & 0 & 0.0 & 8.87 & 31.11 & 58.33 &  \\
20529 & 2 & 2 & 100.0 & 1 & 50.0 & 8.95 & 30.00 & 29.17 &  \\
22480 & 1 & 1 & 100.0 & 0 & 0.0 & 8.00 & 20.00 & 83.33 &  \\
23468 & 1 & 1 & 100.0 & 0 & 0.0 & 9.50 & 33.33 & 75.00 &  \\
23799 & 1 & 1 & 100.0 & 0 & 0.0 & 8.70 & 46.67 & 33.33 &  \\
23824 & 4 & 3 & 75.0 & 0 & 0.0 & 6.50 & 41.67 & 43.75 &  \\
21854 & 8 & 4 & 50.0 & 0 & 0.0 & 5.83 & 56.67 & 47.66 &  \\
22115 & 2 & 0 & 0.0 & 0 & 0.0 & 3.90 & 46.67 & 62.50 &  \\
\bottomrule
\end{longtable}
\endgroup

\clearpage

\subsection{Apendice B: tablas por periodo}

\subsubsection{Anio 2019, sesion OTOÑO}

\begingroup
\scriptsize
\setlength{\tabcolsep}{3pt}
\begin{longtable}{lrrrrrrrrr}
\caption{Tabla completa de profesores para MAT1022, anio 2019, sesion OTOÑO.}\label{tab:prof-period-mat1022-2019-oto-o-especializado}\\
\toprule
Profesor & Total & Obj. & Obj. \% & Bench. & Bench. \% & Cal. & DMU & GA-GB & Rank \\
\midrule
\endfirsthead
\toprule
Profesor & Total & Obj. & Obj. \% & Bench. & Bench. \% & Cal. & DMU & GA-GB & Rank \\
\midrule
\endhead
22471 & 1 & 1 & 100.0 & 0 & 0.0 & 8.30 & 46.67 & 66.67 &  \\
22521 & 1 & 1 & 100.0 & 0 & 0.0 & 8.10 & 46.67 & 75.00 &  \\
\bottomrule
\end{longtable}
\endgroup

\subsubsection{Anio 2019, sesion PRIMAVERA}

\begingroup
\scriptsize
\setlength{\tabcolsep}{3pt}
\begin{longtable}{lrrrrrrrrr}
\caption{Tabla completa de profesores para MAT1022, anio 2019, sesion PRIMAVERA.}\label{tab:prof-period-mat1022-2019-primavera-especializado}\\
\toprule
Profesor & Total & Obj. & Obj. \% & Bench. & Bench. \% & Cal. & DMU & GA-GB & Rank \\
\midrule
\endfirsthead
\toprule
Profesor & Total & Obj. & Obj. \% & Bench. & Bench. \% & Cal. & DMU & GA-GB & Rank \\
\midrule
\endhead
20529 & 2 & 2 & 100.0 & 1 & 50.0 & 8.95 & 30.00 & 29.17 &  \\
22480 & 1 & 1 & 100.0 & 0 & 0.0 & 8.00 & 20.00 & 83.33 &  \\
23153 & 1 & 1 & 100.0 & 0 & 0.0 & 9.70 & 53.33 & 41.67 &  \\
23411 & 1 & 1 & 100.0 & 0 & 0.0 & 9.10 & 60.00 & 58.33 &  \\
23442 & 1 & 1 & 100.0 & 0 & 0.0 & 9.20 & 60.00 & 83.33 &  \\
12516 & 2 & 1 & 50.0 & 0 & 0.0 & 7.60 & 53.33 & 75.00 &  \\
\bottomrule
\end{longtable}
\endgroup

\subsubsection{Anio 2020, sesion OTOÑO}

\begingroup
\scriptsize
\setlength{\tabcolsep}{3pt}
\begin{longtable}{lrrrrrrrrr}
\caption{Tabla completa de profesores para MAT1022, anio 2020, sesion OTOÑO.}\label{tab:prof-period-mat1022-2020-oto-o-especializado}\\
\toprule
Profesor & Total & Obj. & Obj. \% & Bench. & Bench. \% & Cal. & DMU & GA-GB & Rank \\
\midrule
\endfirsthead
\toprule
Profesor & Total & Obj. & Obj. \% & Bench. & Bench. \% & Cal. & DMU & GA-GB & Rank \\
\midrule
\endhead
23162 & 6 & 6 & 100.0 & 1 & 16.7 & 8.93 & 54.44 & 55.56 & 1 \\
23837 & 5 & 5 & 100.0 & 0 & 0.0 & 8.16 & 48.00 & 65.42 & 2 \\
23411 & 4 & 4 & 100.0 & 0 & 0.0 & 8.95 & 51.67 & 75.52 &  \\
23885 & 3 & 3 & 100.0 & 0 & 0.0 & 8.93 & 42.22 & 55.56 &  \\
21262 & 2 & 2 & 100.0 & 0 & 0.0 & 8.45 & 53.33 & 75.00 &  \\
23149 & 2 & 2 & 100.0 & 0 & 0.0 & 7.95 & 46.67 & 54.17 &  \\
23839 & 1 & 1 & 100.0 & 0 & 0.0 & 7.60 & 46.67 & 58.33 &  \\
24097 & 1 & 1 & 100.0 & 0 & 0.0 & 9.40 & 66.67 & 75.00 &  \\
23824 & 4 & 3 & 75.0 & 0 & 0.0 & 6.50 & 41.67 & 43.75 &  \\
21854 & 2 & 0 & 0.0 & 0 & 0.0 & 4.00 & 60.00 & 33.33 &  \\
\bottomrule
\end{longtable}
\endgroup

\subsubsection{Anio 2020, sesion PRIMAVERA}

\begingroup
\scriptsize
\setlength{\tabcolsep}{3pt}
\begin{longtable}{lrrrrrrrrr}
\caption{Tabla completa de profesores para MAT1022, anio 2020, sesion PRIMAVERA.}\label{tab:prof-period-mat1022-2020-primavera-especializado}\\
\toprule
Profesor & Total & Obj. & Obj. \% & Bench. & Bench. \% & Cal. & DMU & GA-GB & Rank \\
\midrule
\endfirsthead
\toprule
Profesor & Total & Obj. & Obj. \% & Bench. & Bench. \% & Cal. & DMU & GA-GB & Rank \\
\midrule
\endhead
23152 & 25 & 25 & 100.0 & 0 & 0.0 & 9.19 & 62.93 & 67.17 & 1 \\
23452 & 25 & 25 & 100.0 & 2 & 8.0 & 9.57 & 48.53 & 59.33 & 2 \\
23885 & 24 & 24 & 100.0 & 1 & 4.2 & 9.28 & 51.67 & 59.64 & 3 \\
23425 & 22 & 22 & 100.0 & 1 & 4.5 & 9.20 & 55.76 & 54.36 & 4 \\
23823 & 22 & 22 & 100.0 & 3 & 13.6 & 8.64 & 50.30 & 56.72 & 5 \\
20792 & 19 & 19 & 100.0 & 0 & 0.0 & 9.04 & 64.56 & 73.25 & 6 \\
22446 & 19 & 19 & 100.0 & 0 & 0.0 & 8.98 & 50.53 & 55.81 & 7 \\
12516 & 17 & 17 & 100.0 & 0 & 0.0 & 9.05 & 52.16 & 59.07 & 8 \\
23411 & 12 & 12 & 100.0 & 0 & 0.0 & 9.43 & 66.11 & 61.28 & 9 \\
23148 & 10 & 10 & 100.0 & 1 & 10.0 & 9.44 & 43.33 & 55.83 & 10 \\
23897 & 49 & 48 & 98.0 & 3 & 6.1 & 9.50 & 53.20 & 58.72 & 11 \\
22114 & 38 & 37 & 97.4 & 2 & 5.3 & 9.19 & 53.68 & 60.96 & 12 \\
23434 & 33 & 32 & 97.0 & 0 & 0.0 & 8.46 & 57.78 & 65.03 & 13 \\
22473 & 25 & 24 & 96.0 & 0 & 0.0 & 9.62 & 53.87 & 58.75 & 14 \\
23442 & 23 & 22 & 95.7 & 1 & 4.3 & 9.07 & 62.32 & 60.51 & 15 \\
24097 & 18 & 17 & 94.4 & 0 & 0.0 & 8.56 & 45.56 & 57.52 & 16 \\
23839 & 32 & 30 & 93.8 & 0 & 0.0 & 8.57 & 60.00 & 55.99 & 17 \\
23836 & 44 & 40 & 90.9 & 1 & 2.3 & 8.42 & 56.97 & 63.30 & 18 \\
23852 & 33 & 30 & 90.9 & 3 & 9.1 & 8.75 & 47.27 & 57.39 & 19 \\
23096 & 22 & 20 & 90.9 & 0 & 0.0 & 8.79 & 58.18 & 65.91 & 20 \\
23837 & 22 & 20 & 90.9 & 0 & 0.0 & 8.50 & 52.12 & 52.37 & 21 \\
22471 & 21 & 19 & 90.5 & 0 & 0.0 & 8.29 & 57.46 & 63.49 & 22 \\
21853 & 20 & 18 & 90.0 & 0 & 0.0 & 8.50 & 61.33 & 64.27 & 23 \\
23835 & 52 & 46 & 88.5 & 0 & 0.0 & 8.67 & 50.51 & 59.98 & 24 \\
23874 & 21 & 18 & 85.7 & 0 & 0.0 & 7.98 & 54.29 & 62.90 & 25 \\
\bottomrule
\end{longtable}
\endgroup

\subsubsection{Anio 2021, sesion OTOÑO}

\begingroup
\scriptsize
\setlength{\tabcolsep}{3pt}
\begin{longtable}{lrrrrrrrrr}
\caption{Tabla completa de profesores para MAT1022, anio 2021, sesion OTOÑO.}\label{tab:prof-period-mat1022-2021-oto-o-especializado}\\
\toprule
Profesor & Total & Obj. & Obj. \% & Bench. & Bench. \% & Cal. & DMU & GA-GB & Rank \\
\midrule
\endfirsthead
\toprule
Profesor & Total & Obj. & Obj. \% & Bench. & Bench. \% & Cal. & DMU & GA-GB & Rank \\
\midrule
\endhead
23425 & 1 & 1 & 100.0 & 1 & 100.0 & 8.90 & 20.00 & 37.50 &  \\
23453 & 1 & 1 & 100.0 & 0 & 0.0 & 7.00 & 20.00 & 43.75 &  \\
23468 & 1 & 1 & 100.0 & 0 & 0.0 & 9.50 & 33.33 & 75.00 &  \\
23837 & 1 & 1 & 100.0 & 0 & 0.0 & 6.70 & 13.33 & 58.33 &  \\
22114 & 1 & 0 & 0.0 & 0 & 0.0 & 2.70 & 20.00 & 56.25 &  \\
23835 & 1 & 0 & 0.0 & 0 & 0.0 & 3.60 & 26.67 & 50.00 &  \\
\bottomrule
\end{longtable}
\endgroup

\subsubsection{Anio 2021, sesion PRIMAVERA}

\begingroup
\scriptsize
\setlength{\tabcolsep}{3pt}
\begin{longtable}{lrrrrrrrrr}
\caption{Tabla completa de profesores para MAT1022, anio 2021, sesion PRIMAVERA.}\label{tab:prof-period-mat1022-2021-primavera-especializado}\\
\toprule
Profesor & Total & Obj. & Obj. \% & Bench. & Bench. \% & Cal. & DMU & GA-GB & Rank \\
\midrule
\endfirsthead
\toprule
Profesor & Total & Obj. & Obj. \% & Bench. & Bench. \% & Cal. & DMU & GA-GB & Rank \\
\midrule
\endhead
22471 & 1 & 1 & 100.0 & 0 & 0.0 & 7.20 & 60.00 & 41.67 &  \\
23799 & 1 & 1 & 100.0 & 0 & 0.0 & 8.70 & 46.67 & 33.33 &  \\
24085 & 1 & 1 & 100.0 & 0 & 0.0 & 9.00 & 33.33 & 50.00 &  \\
20792 & 1 & 0 & 0.0 & 0 & 0.0 & 2.20 & 26.67 & 91.67 &  \\
22521 & 1 & 0 & 0.0 & 0 & 0.0 & 2.10 & 53.33 & 33.33 &  \\
\bottomrule
\end{longtable}
\endgroup

\subsubsection{Anio 2022, sesion OTOÑO}

\begingroup
\scriptsize
\setlength{\tabcolsep}{3pt}
\begin{longtable}{lrrrrrrrrr}
\caption{Tabla completa de profesores para MAT1022, anio 2022, sesion OTOÑO.}\label{tab:prof-period-mat1022-2022-oto-o-especializado}\\
\toprule
Profesor & Total & Obj. & Obj. \% & Bench. & Bench. \% & Cal. & DMU & GA-GB & Rank \\
\midrule
\endfirsthead
\toprule
Profesor & Total & Obj. & Obj. \% & Bench. & Bench. \% & Cal. & DMU & GA-GB & Rank \\
\midrule
\endhead
23452 & 7 & 5 & 71.4 & 0 & 0.0 & 8.19 & 75.24 & 57.44 & 1 \\
23897 & 5 & 2 & 40.0 & 0 & 0.0 & 5.72 & 50.67 & 56.67 & 2 \\
24827 & 2 & 2 & 100.0 & 0 & 0.0 & 8.90 & 56.67 & 87.50 &  \\
23874 & 4 & 3 & 75.0 & 0 & 0.0 & 7.05 & 30.00 & 68.23 &  \\
21192 & 2 & 1 & 50.0 & 0 & 0.0 & 10.00 & 73.33 & 41.67 &  \\
23885 & 2 & 1 & 50.0 & 0 & 0.0 & 5.45 & 70.00 & 79.17 &  \\
22471 & 2 & 0 & 0.0 & 0 & 0.0 & 0.35 & 63.33 & 70.83 &  \\
23411 & 1 & 0 & 0.0 & 0 & 0.0 & 5.10 & 73.33 & 58.33 &  \\
\bottomrule
\end{longtable}
\endgroup

\subsubsection{Anio 2022, sesion PRIMAVERA}

\begingroup
\scriptsize
\setlength{\tabcolsep}{3pt}
\begin{longtable}{lrrrrrrrrr}
\caption{Tabla completa de profesores para MAT1022, anio 2022, sesion PRIMAVERA.}\label{tab:prof-period-mat1022-2022-primavera-especializado}\\
\toprule
Profesor & Total & Obj. & Obj. \% & Bench. & Bench. \% & Cal. & DMU & GA-GB & Rank \\
\midrule
\endfirsthead
\toprule
Profesor & Total & Obj. & Obj. \% & Bench. & Bench. \% & Cal. & DMU & GA-GB & Rank \\
\midrule
\endhead
24085 & 42 & 42 & 100.0 & 0 & 0.0 & 9.84 & 63.17 & 73.36 & 1 \\
23425 & 40 & 40 & 100.0 & 1 & 2.5 & 9.36 & 63.83 & 70.78 & 2 \\
21853 & 20 & 20 & 100.0 & 0 & 0.0 & 9.29 & 72.67 & 75.42 & 3 \\
23452 & 20 & 20 & 100.0 & 0 & 0.0 & 9.52 & 79.67 & 74.58 & 4 \\
22521 & 14 & 14 & 100.0 & 0 & 0.0 & 8.74 & 70.00 & 73.66 & 5 \\
22446 & 44 & 43 & 97.7 & 1 & 2.3 & 9.59 & 59.55 & 70.83 & 6 \\
23897 & 40 & 39 & 97.5 & 1 & 2.5 & 9.59 & 72.67 & 71.20 & 7 \\
23411 & 60 & 57 & 95.0 & 0 & 0.0 & 8.97 & 62.56 & 72.12 & 8 \\
23453 & 33 & 31 & 93.9 & 1 & 3.0 & 8.24 & 57.98 & 71.02 & 9 \\
23836 & 12 & 11 & 91.7 & 1 & 8.3 & 8.88 & 66.67 & 66.15 & 10 \\
23148 & 11 & 10 & 90.9 & 0 & 0.0 & 8.05 & 53.94 & 65.15 & 11 \\
21192 & 43 & 39 & 90.7 & 0 & 0.0 & 9.07 & 68.99 & 66.91 & 12 \\
18269 & 36 & 32 & 88.9 & 1 & 2.8 & 8.62 & 71.11 & 75.64 & 13 \\
23835 & 20 & 17 & 85.0 & 0 & 0.0 & 7.61 & 78.00 & 78.75 & 14 \\
21779 & 9 & 6 & 66.7 & 0 & 0.0 & 6.67 & 73.33 & 71.76 &  \\
\bottomrule
\end{longtable}
\endgroup

\subsubsection{Anio 2023, sesion OTOÑO}

\begingroup
\scriptsize
\setlength{\tabcolsep}{3pt}
\begin{longtable}{lrrrrrrrrr}
\caption{Tabla completa de profesores para MAT1022, anio 2023, sesion OTOÑO.}\label{tab:prof-period-mat1022-2023-oto-o-especializado}\\
\toprule
Profesor & Total & Obj. & Obj. \% & Bench. & Bench. \% & Cal. & DMU & GA-GB & Rank \\
\midrule
\endfirsthead
\toprule
Profesor & Total & Obj. & Obj. \% & Bench. & Bench. \% & Cal. & DMU & GA-GB & Rank \\
\midrule
\endhead
23852 & 7 & 7 & 100.0 & 0 & 0.0 & 8.96 & 69.52 & 66.37 & 1 \\
23452 & 8 & 7 & 87.5 & 0 & 0.0 & 8.49 & 55.83 & 55.99 & 2 \\
21192 & 5 & 4 & 80.0 & 0 & 0.0 & 8.96 & 54.67 & 48.33 & 3 \\
22114 & 4 & 4 & 100.0 & 0 & 0.0 & 8.70 & 53.33 & 48.96 &  \\
23887 & 3 & 3 & 100.0 & 0 & 0.0 & 8.87 & 31.11 & 58.33 &  \\
24086 & 2 & 1 & 50.0 & 0 & 0.0 & 5.25 & 20.00 & 25.00 &  \\
23152 & 3 & 0 & 0.0 & 0 & 0.0 & 2.67 & 42.22 & 61.11 &  \\
22115 & 2 & 0 & 0.0 & 0 & 0.0 & 3.90 & 46.67 & 62.50 &  \\
23823 & 2 & 0 & 0.0 & 0 & 0.0 & 2.15 & 73.33 & 66.67 &  \\
\bottomrule
\end{longtable}
\endgroup

\subsubsection{Anio 2023, sesion PRIMAVERA}

\begingroup
\scriptsize
\setlength{\tabcolsep}{3pt}
\begin{longtable}{lrrrrrrrrr}
\caption{Tabla completa de profesores para MAT1022, anio 2023, sesion PRIMAVERA.}\label{tab:prof-period-mat1022-2023-primavera-especializado}\\
\toprule
Profesor & Total & Obj. & Obj. \% & Bench. & Bench. \% & Cal. & DMU & GA-GB & Rank \\
\midrule
\endfirsthead
\toprule
Profesor & Total & Obj. & Obj. \% & Bench. & Bench. \% & Cal. & DMU & GA-GB & Rank \\
\midrule
\endhead
23885 & 42 & 42 & 100.0 & 1 & 2.4 & 9.30 & 67.14 & 71.23 & 1 \\
23897 & 25 & 25 & 100.0 & 0 & 0.0 & 9.44 & 79.47 & 77.75 & 2 \\
23852 & 22 & 22 & 100.0 & 0 & 0.0 & 9.34 & 78.79 & 74.05 & 3 \\
20792 & 19 & 19 & 100.0 & 1 & 5.3 & 8.31 & 56.14 & 71.60 & 4 \\
23148 & 19 & 19 & 100.0 & 0 & 0.0 & 8.71 & 59.30 & 70.83 & 5 \\
22473 & 26 & 25 & 96.2 & 0 & 0.0 & 9.03 & 69.23 & 70.99 & 6 \\
23149 & 26 & 25 & 96.2 & 0 & 0.0 & 9.22 & 66.41 & 73.48 & 7 \\
24827 & 24 & 23 & 95.8 & 0 & 0.0 & 9.32 & 75.00 & 70.75 & 8 \\
24085 & 18 & 16 & 88.9 & 2 & 11.1 & 9.96 & 52.59 & 69.79 & 9 \\
22446 & 17 & 15 & 88.2 & 0 & 0.0 & 8.66 & 72.16 & 69.73 & 10 \\
22521 & 14 & 11 & 78.6 & 0 & 0.0 & 8.17 & 66.19 & 79.91 & 11 \\
23836 & 16 & 12 & 75.0 & 0 & 0.0 & 7.21 & 65.83 & 67.32 & 12 \\
21779 & 23 & 17 & 73.9 & 0 & 0.0 & 8.10 & 65.22 & 64.67 & 13 \\
23411 & 6 & 6 & 100.0 & 0 & 0.0 & 9.43 & 77.78 & 72.22 &  \\
\bottomrule
\end{longtable}
\endgroup

\subsubsection{Anio 2024, sesion OTOÑO}

\begingroup
\scriptsize
\setlength{\tabcolsep}{3pt}
\begin{longtable}{lrrrrrrrrr}
\caption{Tabla completa de profesores para MAT1022, anio 2024, sesion OTOÑO.}\label{tab:prof-period-mat1022-2024-oto-o-especializado}\\
\toprule
Profesor & Total & Obj. & Obj. \% & Bench. & Bench. \% & Cal. & DMU & GA-GB & Rank \\
\midrule
\endfirsthead
\toprule
Profesor & Total & Obj. & Obj. \% & Bench. & Bench. \% & Cal. & DMU & GA-GB & Rank \\
\midrule
\endhead
21192 & 10 & 10 & 100.0 & 2 & 20.0 & 9.37 & 58.00 & 51.67 & 1 \\
25503 & 2 & 2 & 100.0 & 0 & 0.0 & 8.50 & 50.00 & 58.33 &  \\
23452 & 1 & 1 & 100.0 & 0 & 0.0 & 9.10 & 66.67 & 41.67 &  \\
22473 & 5 & 4 & 80.0 & 0 & 0.0 & 7.48 & 65.33 & 53.33 &  \\
23148 & 5 & 4 & 80.0 & 0 & 0.0 & 7.90 & 60.00 & 49.58 &  \\
21854 & 6 & 4 & 66.7 & 0 & 0.0 & 6.43 & 55.56 & 52.43 &  \\
24086 & 5 & 2 & 40.0 & 0 & 0.0 & 6.38 & 53.33 & 54.17 &  \\
23152 & 3 & 0 & 0.0 & 0 & 0.0 & 3.07 & 80.00 & 59.72 &  \\
22114 & 1 & 0 & 0.0 & 0 & 0.0 & 1.10 & 53.33 & 50.00 &  \\
\bottomrule
\end{longtable}
\endgroup

\subsubsection{Anio 2024, sesion PRIMAVERA}

\begingroup
\scriptsize
\setlength{\tabcolsep}{3pt}
\begin{longtable}{lrrrrrrrrr}
\caption{Tabla completa de profesores para MAT1022, anio 2024, sesion PRIMAVERA.}\label{tab:prof-period-mat1022-2024-primavera-especializado}\\
\toprule
Profesor & Total & Obj. & Obj. \% & Bench. & Bench. \% & Cal. & DMU & GA-GB & Rank \\
\midrule
\endfirsthead
\toprule
Profesor & Total & Obj. & Obj. \% & Bench. & Bench. \% & Cal. & DMU & GA-GB & Rank \\
\midrule
\endhead
23452 & 29 & 29 & 100.0 & 0 & 0.0 & 9.39 & 79.31 & 72.05 & 1 \\
24097 & 29 & 29 & 100.0 & 0 & 0.0 & 9.53 & 80.92 & 60.70 & 2 \\
23885 & 28 & 28 & 100.0 & 0 & 0.0 & 9.15 & 66.43 & 70.31 & 3 \\
23823 & 17 & 17 & 100.0 & 0 & 0.0 & 9.31 & 76.47 & 60.78 & 4 \\
23978 & 15 & 15 & 100.0 & 0 & 0.0 & 9.20 & 68.00 & 65.69 & 5 \\
21779 & 34 & 33 & 97.1 & 1 & 2.9 & 8.89 & 70.00 & 71.45 & 6 \\
23836 & 55 & 53 & 96.4 & 2 & 3.6 & 9.55 & 63.88 & 62.65 & 7 \\
23852 & 55 & 53 & 96.4 & 0 & 0.0 & 9.19 & 68.36 & 65.87 & 8 \\
24086 & 26 & 25 & 96.2 & 0 & 0.0 & 8.89 & 71.28 & 74.68 & 9 \\
23152 & 22 & 21 & 95.5 & 0 & 0.0 & 9.00 & 74.24 & 71.40 & 10 \\
23153 & 21 & 20 & 95.2 & 0 & 0.0 & 8.90 & 64.44 & 64.98 & 11 \\
21262 & 19 & 18 & 94.7 & 0 & 0.0 & 8.25 & 74.39 & 69.08 & 12 \\
23148 & 19 & 18 & 94.7 & 0 & 0.0 & 8.77 & 71.93 & 69.08 & 13 \\
24085 & 46 & 43 & 93.5 & 1 & 2.2 & 9.98 & 63.62 & 65.31 & 14 \\
23411 & 27 & 25 & 92.6 & 0 & 0.0 & 9.03 & 74.32 & 66.28 & 15 \\
22446 & 13 & 12 & 92.3 & 0 & 0.0 & 8.28 & 64.62 & 69.39 & 16 \\
23839 & 22 & 20 & 90.9 & 0 & 0.0 & 8.37 & 58.18 & 65.72 & 17 \\
22471 & 14 & 12 & 85.7 & 0 & 0.0 & 8.00 & 63.33 & 67.26 & 18 \\
\bottomrule
\end{longtable}
\endgroup

\subsubsection{Anio 2025, sesion PRIMAVERA}

\begingroup
\scriptsize
\setlength{\tabcolsep}{3pt}
\begin{longtable}{lrrrrrrrrr}
\caption{Tabla completa de profesores para MAT1022, anio 2025, sesion PRIMAVERA.}\label{tab:prof-period-mat1022-2025-primavera-especializado}\\
\toprule
Profesor & Total & Obj. & Obj. \% & Bench. & Bench. \% & Cal. & DMU & GA-GB & Rank \\
\midrule
\endfirsthead
\toprule
Profesor & Total & Obj. & Obj. \% & Bench. & Bench. \% & Cal. & DMU & GA-GB & Rank \\
\midrule
\endhead
23836 & 39 & 39 & 100.0 & 0 & 0.0 & 9.58 & 78.12 & 74.41 & 1 \\
23885 & 26 & 26 & 100.0 & 0 & 0.0 & 9.35 & 76.67 & 73.40 & 2 \\
24085 & 26 & 26 & 100.0 & 0 & 0.0 & 9.55 & 63.85 & 65.38 & 3 \\
25503 & 20 & 20 & 100.0 & 0 & 0.0 & 9.80 & 69.00 & 68.33 & 4 \\
24097 & 16 & 16 & 100.0 & 0 & 0.0 & 9.68 & 88.75 & 74.87 & 5 \\
23153 & 15 & 15 & 100.0 & 0 & 0.0 & 9.69 & 78.67 & 70.42 & 6 \\
23411 & 10 & 10 & 100.0 & 0 & 0.0 & 9.57 & 81.33 & 74.17 & 7 \\
25531 & 10 & 10 & 100.0 & 0 & 0.0 & 9.87 & 70.67 & 71.67 & 8 \\
23852 & 41 & 40 & 97.6 & 1 & 2.4 & 9.06 & 75.28 & 71.24 & 9 \\
24086 & 35 & 34 & 97.1 & 0 & 0.0 & 9.37 & 72.95 & 76.19 & 10 \\
23978 & 34 & 33 & 97.1 & 0 & 0.0 & 9.42 & 74.71 & 66.91 & 11 \\
22446 & 19 & 18 & 94.7 & 0 & 0.0 & 9.28 & 81.05 & 58.99 & 12 \\
20792 & 17 & 16 & 94.1 & 0 & 0.0 & 9.52 & 85.10 & 79.41 & 13 \\
21779 & 31 & 28 & 90.3 & 0 & 0.0 & 8.53 & 71.40 & 72.31 & 14 \\
23452 & 20 & 18 & 90.0 & 0 & 0.0 & 9.46 & 76.00 & 63.54 & 15 \\
25843 & 23 & 20 & 87.0 & 0 & 0.0 & 8.01 & 71.30 & 63.59 & 16 \\
21262 & 8 & 8 & 100.0 & 0 & 0.0 & 8.21 & 68.33 & 64.58 &  \\
23152 & 8 & 8 & 100.0 & 0 & 0.0 & 8.49 & 62.50 & 69.79 &  \\
23823 & 3 & 3 & 100.0 & 0 & 0.0 & 9.77 & 64.44 & 55.56 &  \\
\bottomrule
\end{longtable}
\endgroup}{\IfFileExists{apendice_tablas_profesores_MAT1022.tex}{\clearpage

\appendix

\section{Tablas completas de profesores para MAT1022}

Este apendice se genera automaticamente desde \code{data/datos\_procesados/professor\_appendix\_all\_years} y \code{data/datos\_procesados/professor\_appendix\_by\_period}. El desglose temporal usa exactamente las combinaciones observadas de \code{anio} y \code{CLAVESESION}.

CÁLCULO I

\subsection{Apendice A: tabla global de profesores}

\begingroup
\scriptsize
\setlength{\tabcolsep}{3pt}
\begin{longtable}{lrrrrrrrrr}
\caption{Tabla completa de profesores para MAT1022, todos los anios.}\label{tab:prof-global-mat1022-especializado}\\
\toprule
Profesor & Total & Obj. & Obj. \% & Bench. & Bench. \% & Cal. & DMU & GA-GB & Rank \\
\midrule
\endfirsthead
\toprule
Profesor & Total & Obj. & Obj. \% & Bench. & Bench. \% & Cal. & DMU & GA-GB & Rank \\
\midrule
\endhead
23425 & 63 & 63 & 100.0 & 3 & 4.8 & 9.30 & 60.32 & 64.52 & 1 \\
25503 & 22 & 22 & 100.0 & 0 & 0.0 & 9.69 & 67.27 & 67.42 & 2 \\
25531 & 10 & 10 & 100.0 & 0 & 0.0 & 9.87 & 70.67 & 71.67 & 3 \\
23885 & 125 & 124 & 99.2 & 2 & 1.6 & 9.20 & 65.44 & 69.00 & 4 \\
24097 & 64 & 63 & 98.4 & 0 & 0.0 & 9.29 & 72.71 & 63.57 & 5 \\
23978 & 49 & 48 & 98.0 & 0 & 0.0 & 9.35 & 72.65 & 66.54 & 6 \\
23153 & 37 & 36 & 97.3 & 0 & 0.0 & 9.24 & 69.91 & 66.55 & 7 \\
23434 & 33 & 32 & 97.0 & 0 & 0.0 & 8.46 & 57.78 & 65.03 & 8 \\
21262 & 29 & 28 & 96.6 & 0 & 0.0 & 8.26 & 71.26 & 68.25 & 9 \\
20792 & 56 & 54 & 96.4 & 1 & 1.8 & 8.82 & 67.26 & 74.89 & 10 \\
23149 & 28 & 27 & 96.4 & 0 & 0.0 & 9.13 & 65.00 & 72.10 & 11 \\
24085 & 133 & 128 & 96.2 & 3 & 2.3 & 9.84 & 61.80 & 68.36 & 12 \\
23852 & 158 & 152 & 96.2 & 4 & 2.5 & 9.07 & 67.26 & 66.65 & 13 \\
24827 & 26 & 25 & 96.2 & 0 & 0.0 & 9.29 & 73.59 & 72.04 & 14 \\
23442 & 24 & 23 & 95.8 & 1 & 4.2 & 9.07 & 62.22 & 61.46 & 15 \\
23897 & 119 & 114 & 95.8 & 4 & 3.4 & 9.36 & 65.15 & 66.82 & 16 \\
22446 & 112 & 107 & 95.5 & 1 & 0.9 & 9.14 & 64.17 & 65.94 & 17 \\
23452 & 110 & 105 & 95.5 & 2 & 1.8 & 9.32 & 69.70 & 65.70 & 18 \\
23823 & 44 & 42 & 95.5 & 3 & 6.8 & 8.68 & 62.42 & 58.66 & 19 \\
23148 & 64 & 61 & 95.3 & 1 & 1.6 & 8.66 & 59.69 & 65.33 & 20 \\
23411 & 121 & 115 & 95.0 & 0 & 0.0 & 9.07 & 67.55 & 69.80 & 21 \\
21853 & 40 & 38 & 95.0 & 0 & 0.0 & 8.90 & 67.00 & 69.84 & 22 \\
12516 & 19 & 18 & 94.7 & 0 & 0.0 & 8.90 & 52.28 & 60.75 & 23 \\
22473 & 56 & 53 & 94.6 & 0 & 0.0 & 9.16 & 62.02 & 63.95 & 24 \\
23453 & 34 & 32 & 94.1 & 1 & 2.9 & 8.20 & 56.86 & 70.22 & 25 \\
23836 & 166 & 155 & 93.4 & 4 & 2.4 & 8.98 & 65.78 & 66.29 & 26 \\
22114 & 44 & 41 & 93.2 & 2 & 4.5 & 8.81 & 52.88 & 59.52 & 27 \\
23837 & 28 & 26 & 92.9 & 0 & 0.0 & 8.38 & 50.00 & 54.91 & 28 \\
23839 & 55 & 51 & 92.7 & 0 & 0.0 & 8.47 & 59.03 & 59.92 & 29 \\
24086 & 68 & 62 & 91.2 & 0 & 0.0 & 8.85 & 69.31 & 72.49 & 30 \\
23096 & 22 & 20 & 90.9 & 0 & 0.0 & 8.79 & 58.18 & 65.91 & 31 \\
21192 & 60 & 54 & 90.0 & 2 & 3.3 & 9.14 & 66.11 & 61.98 & 32 \\
18269 & 36 & 32 & 88.9 & 1 & 2.8 & 8.62 & 71.11 & 75.64 & 33 \\
23152 & 61 & 54 & 88.5 & 0 & 0.0 & 8.41 & 66.78 & 68.37 & 34 \\
25843 & 23 & 20 & 87.0 & 0 & 0.0 & 8.01 & 71.30 & 63.59 & 35 \\
22521 & 30 & 26 & 86.7 & 0 & 0.0 & 8.23 & 66.89 & 75.28 & 36 \\
21779 & 97 & 84 & 86.6 & 1 & 1.0 & 8.38 & 69.62 & 70.15 & 37 \\
23835 & 73 & 63 & 86.3 & 0 & 0.0 & 8.31 & 57.72 & 64.98 & 38 \\
22471 & 39 & 33 & 84.6 & 0 & 0.0 & 7.75 & 59.66 & 64.74 & 39 \\
23874 & 25 & 21 & 84.0 & 0 & 0.0 & 7.83 & 50.40 & 63.75 & 40 \\
23162 & 6 & 6 & 100.0 & 1 & 16.7 & 8.93 & 54.44 & 55.56 &  \\
23887 & 3 & 3 & 100.0 & 0 & 0.0 & 8.87 & 31.11 & 58.33 &  \\
20529 & 2 & 2 & 100.0 & 1 & 50.0 & 8.95 & 30.00 & 29.17 &  \\
22480 & 1 & 1 & 100.0 & 0 & 0.0 & 8.00 & 20.00 & 83.33 &  \\
23468 & 1 & 1 & 100.0 & 0 & 0.0 & 9.50 & 33.33 & 75.00 &  \\
23799 & 1 & 1 & 100.0 & 0 & 0.0 & 8.70 & 46.67 & 33.33 &  \\
23824 & 4 & 3 & 75.0 & 0 & 0.0 & 6.50 & 41.67 & 43.75 &  \\
21854 & 8 & 4 & 50.0 & 0 & 0.0 & 5.83 & 56.67 & 47.66 &  \\
22115 & 2 & 0 & 0.0 & 0 & 0.0 & 3.90 & 46.67 & 62.50 &  \\
\bottomrule
\end{longtable}
\endgroup

\clearpage

\subsection{Apendice B: tablas por periodo}

\subsubsection{Anio 2019, sesion OTOÑO}

\begingroup
\scriptsize
\setlength{\tabcolsep}{3pt}
\begin{longtable}{lrrrrrrrrr}
\caption{Tabla completa de profesores para MAT1022, anio 2019, sesion OTOÑO.}\label{tab:prof-period-mat1022-2019-oto-o-especializado}\\
\toprule
Profesor & Total & Obj. & Obj. \% & Bench. & Bench. \% & Cal. & DMU & GA-GB & Rank \\
\midrule
\endfirsthead
\toprule
Profesor & Total & Obj. & Obj. \% & Bench. & Bench. \% & Cal. & DMU & GA-GB & Rank \\
\midrule
\endhead
22471 & 1 & 1 & 100.0 & 0 & 0.0 & 8.30 & 46.67 & 66.67 &  \\
22521 & 1 & 1 & 100.0 & 0 & 0.0 & 8.10 & 46.67 & 75.00 &  \\
\bottomrule
\end{longtable}
\endgroup

\subsubsection{Anio 2019, sesion PRIMAVERA}

\begingroup
\scriptsize
\setlength{\tabcolsep}{3pt}
\begin{longtable}{lrrrrrrrrr}
\caption{Tabla completa de profesores para MAT1022, anio 2019, sesion PRIMAVERA.}\label{tab:prof-period-mat1022-2019-primavera-especializado}\\
\toprule
Profesor & Total & Obj. & Obj. \% & Bench. & Bench. \% & Cal. & DMU & GA-GB & Rank \\
\midrule
\endfirsthead
\toprule
Profesor & Total & Obj. & Obj. \% & Bench. & Bench. \% & Cal. & DMU & GA-GB & Rank \\
\midrule
\endhead
20529 & 2 & 2 & 100.0 & 1 & 50.0 & 8.95 & 30.00 & 29.17 &  \\
22480 & 1 & 1 & 100.0 & 0 & 0.0 & 8.00 & 20.00 & 83.33 &  \\
23153 & 1 & 1 & 100.0 & 0 & 0.0 & 9.70 & 53.33 & 41.67 &  \\
23411 & 1 & 1 & 100.0 & 0 & 0.0 & 9.10 & 60.00 & 58.33 &  \\
23442 & 1 & 1 & 100.0 & 0 & 0.0 & 9.20 & 60.00 & 83.33 &  \\
12516 & 2 & 1 & 50.0 & 0 & 0.0 & 7.60 & 53.33 & 75.00 &  \\
\bottomrule
\end{longtable}
\endgroup

\subsubsection{Anio 2020, sesion OTOÑO}

\begingroup
\scriptsize
\setlength{\tabcolsep}{3pt}
\begin{longtable}{lrrrrrrrrr}
\caption{Tabla completa de profesores para MAT1022, anio 2020, sesion OTOÑO.}\label{tab:prof-period-mat1022-2020-oto-o-especializado}\\
\toprule
Profesor & Total & Obj. & Obj. \% & Bench. & Bench. \% & Cal. & DMU & GA-GB & Rank \\
\midrule
\endfirsthead
\toprule
Profesor & Total & Obj. & Obj. \% & Bench. & Bench. \% & Cal. & DMU & GA-GB & Rank \\
\midrule
\endhead
23162 & 6 & 6 & 100.0 & 1 & 16.7 & 8.93 & 54.44 & 55.56 & 1 \\
23837 & 5 & 5 & 100.0 & 0 & 0.0 & 8.16 & 48.00 & 65.42 & 2 \\
23411 & 4 & 4 & 100.0 & 0 & 0.0 & 8.95 & 51.67 & 75.52 &  \\
23885 & 3 & 3 & 100.0 & 0 & 0.0 & 8.93 & 42.22 & 55.56 &  \\
21262 & 2 & 2 & 100.0 & 0 & 0.0 & 8.45 & 53.33 & 75.00 &  \\
23149 & 2 & 2 & 100.0 & 0 & 0.0 & 7.95 & 46.67 & 54.17 &  \\
23839 & 1 & 1 & 100.0 & 0 & 0.0 & 7.60 & 46.67 & 58.33 &  \\
24097 & 1 & 1 & 100.0 & 0 & 0.0 & 9.40 & 66.67 & 75.00 &  \\
23824 & 4 & 3 & 75.0 & 0 & 0.0 & 6.50 & 41.67 & 43.75 &  \\
21854 & 2 & 0 & 0.0 & 0 & 0.0 & 4.00 & 60.00 & 33.33 &  \\
\bottomrule
\end{longtable}
\endgroup

\subsubsection{Anio 2020, sesion PRIMAVERA}

\begingroup
\scriptsize
\setlength{\tabcolsep}{3pt}
\begin{longtable}{lrrrrrrrrr}
\caption{Tabla completa de profesores para MAT1022, anio 2020, sesion PRIMAVERA.}\label{tab:prof-period-mat1022-2020-primavera-especializado}\\
\toprule
Profesor & Total & Obj. & Obj. \% & Bench. & Bench. \% & Cal. & DMU & GA-GB & Rank \\
\midrule
\endfirsthead
\toprule
Profesor & Total & Obj. & Obj. \% & Bench. & Bench. \% & Cal. & DMU & GA-GB & Rank \\
\midrule
\endhead
23152 & 25 & 25 & 100.0 & 0 & 0.0 & 9.19 & 62.93 & 67.17 & 1 \\
23452 & 25 & 25 & 100.0 & 2 & 8.0 & 9.57 & 48.53 & 59.33 & 2 \\
23885 & 24 & 24 & 100.0 & 1 & 4.2 & 9.28 & 51.67 & 59.64 & 3 \\
23425 & 22 & 22 & 100.0 & 1 & 4.5 & 9.20 & 55.76 & 54.36 & 4 \\
23823 & 22 & 22 & 100.0 & 3 & 13.6 & 8.64 & 50.30 & 56.72 & 5 \\
20792 & 19 & 19 & 100.0 & 0 & 0.0 & 9.04 & 64.56 & 73.25 & 6 \\
22446 & 19 & 19 & 100.0 & 0 & 0.0 & 8.98 & 50.53 & 55.81 & 7 \\
12516 & 17 & 17 & 100.0 & 0 & 0.0 & 9.05 & 52.16 & 59.07 & 8 \\
23411 & 12 & 12 & 100.0 & 0 & 0.0 & 9.43 & 66.11 & 61.28 & 9 \\
23148 & 10 & 10 & 100.0 & 1 & 10.0 & 9.44 & 43.33 & 55.83 & 10 \\
23897 & 49 & 48 & 98.0 & 3 & 6.1 & 9.50 & 53.20 & 58.72 & 11 \\
22114 & 38 & 37 & 97.4 & 2 & 5.3 & 9.19 & 53.68 & 60.96 & 12 \\
23434 & 33 & 32 & 97.0 & 0 & 0.0 & 8.46 & 57.78 & 65.03 & 13 \\
22473 & 25 & 24 & 96.0 & 0 & 0.0 & 9.62 & 53.87 & 58.75 & 14 \\
23442 & 23 & 22 & 95.7 & 1 & 4.3 & 9.07 & 62.32 & 60.51 & 15 \\
24097 & 18 & 17 & 94.4 & 0 & 0.0 & 8.56 & 45.56 & 57.52 & 16 \\
23839 & 32 & 30 & 93.8 & 0 & 0.0 & 8.57 & 60.00 & 55.99 & 17 \\
23836 & 44 & 40 & 90.9 & 1 & 2.3 & 8.42 & 56.97 & 63.30 & 18 \\
23852 & 33 & 30 & 90.9 & 3 & 9.1 & 8.75 & 47.27 & 57.39 & 19 \\
23096 & 22 & 20 & 90.9 & 0 & 0.0 & 8.79 & 58.18 & 65.91 & 20 \\
23837 & 22 & 20 & 90.9 & 0 & 0.0 & 8.50 & 52.12 & 52.37 & 21 \\
22471 & 21 & 19 & 90.5 & 0 & 0.0 & 8.29 & 57.46 & 63.49 & 22 \\
21853 & 20 & 18 & 90.0 & 0 & 0.0 & 8.50 & 61.33 & 64.27 & 23 \\
23835 & 52 & 46 & 88.5 & 0 & 0.0 & 8.67 & 50.51 & 59.98 & 24 \\
23874 & 21 & 18 & 85.7 & 0 & 0.0 & 7.98 & 54.29 & 62.90 & 25 \\
\bottomrule
\end{longtable}
\endgroup

\subsubsection{Anio 2021, sesion OTOÑO}

\begingroup
\scriptsize
\setlength{\tabcolsep}{3pt}
\begin{longtable}{lrrrrrrrrr}
\caption{Tabla completa de profesores para MAT1022, anio 2021, sesion OTOÑO.}\label{tab:prof-period-mat1022-2021-oto-o-especializado}\\
\toprule
Profesor & Total & Obj. & Obj. \% & Bench. & Bench. \% & Cal. & DMU & GA-GB & Rank \\
\midrule
\endfirsthead
\toprule
Profesor & Total & Obj. & Obj. \% & Bench. & Bench. \% & Cal. & DMU & GA-GB & Rank \\
\midrule
\endhead
23425 & 1 & 1 & 100.0 & 1 & 100.0 & 8.90 & 20.00 & 37.50 &  \\
23453 & 1 & 1 & 100.0 & 0 & 0.0 & 7.00 & 20.00 & 43.75 &  \\
23468 & 1 & 1 & 100.0 & 0 & 0.0 & 9.50 & 33.33 & 75.00 &  \\
23837 & 1 & 1 & 100.0 & 0 & 0.0 & 6.70 & 13.33 & 58.33 &  \\
22114 & 1 & 0 & 0.0 & 0 & 0.0 & 2.70 & 20.00 & 56.25 &  \\
23835 & 1 & 0 & 0.0 & 0 & 0.0 & 3.60 & 26.67 & 50.00 &  \\
\bottomrule
\end{longtable}
\endgroup

\subsubsection{Anio 2021, sesion PRIMAVERA}

\begingroup
\scriptsize
\setlength{\tabcolsep}{3pt}
\begin{longtable}{lrrrrrrrrr}
\caption{Tabla completa de profesores para MAT1022, anio 2021, sesion PRIMAVERA.}\label{tab:prof-period-mat1022-2021-primavera-especializado}\\
\toprule
Profesor & Total & Obj. & Obj. \% & Bench. & Bench. \% & Cal. & DMU & GA-GB & Rank \\
\midrule
\endfirsthead
\toprule
Profesor & Total & Obj. & Obj. \% & Bench. & Bench. \% & Cal. & DMU & GA-GB & Rank \\
\midrule
\endhead
22471 & 1 & 1 & 100.0 & 0 & 0.0 & 7.20 & 60.00 & 41.67 &  \\
23799 & 1 & 1 & 100.0 & 0 & 0.0 & 8.70 & 46.67 & 33.33 &  \\
24085 & 1 & 1 & 100.0 & 0 & 0.0 & 9.00 & 33.33 & 50.00 &  \\
20792 & 1 & 0 & 0.0 & 0 & 0.0 & 2.20 & 26.67 & 91.67 &  \\
22521 & 1 & 0 & 0.0 & 0 & 0.0 & 2.10 & 53.33 & 33.33 &  \\
\bottomrule
\end{longtable}
\endgroup

\subsubsection{Anio 2022, sesion OTOÑO}

\begingroup
\scriptsize
\setlength{\tabcolsep}{3pt}
\begin{longtable}{lrrrrrrrrr}
\caption{Tabla completa de profesores para MAT1022, anio 2022, sesion OTOÑO.}\label{tab:prof-period-mat1022-2022-oto-o-especializado}\\
\toprule
Profesor & Total & Obj. & Obj. \% & Bench. & Bench. \% & Cal. & DMU & GA-GB & Rank \\
\midrule
\endfirsthead
\toprule
Profesor & Total & Obj. & Obj. \% & Bench. & Bench. \% & Cal. & DMU & GA-GB & Rank \\
\midrule
\endhead
23452 & 7 & 5 & 71.4 & 0 & 0.0 & 8.19 & 75.24 & 57.44 & 1 \\
23897 & 5 & 2 & 40.0 & 0 & 0.0 & 5.72 & 50.67 & 56.67 & 2 \\
24827 & 2 & 2 & 100.0 & 0 & 0.0 & 8.90 & 56.67 & 87.50 &  \\
23874 & 4 & 3 & 75.0 & 0 & 0.0 & 7.05 & 30.00 & 68.23 &  \\
21192 & 2 & 1 & 50.0 & 0 & 0.0 & 10.00 & 73.33 & 41.67 &  \\
23885 & 2 & 1 & 50.0 & 0 & 0.0 & 5.45 & 70.00 & 79.17 &  \\
22471 & 2 & 0 & 0.0 & 0 & 0.0 & 0.35 & 63.33 & 70.83 &  \\
23411 & 1 & 0 & 0.0 & 0 & 0.0 & 5.10 & 73.33 & 58.33 &  \\
\bottomrule
\end{longtable}
\endgroup

\subsubsection{Anio 2022, sesion PRIMAVERA}

\begingroup
\scriptsize
\setlength{\tabcolsep}{3pt}
\begin{longtable}{lrrrrrrrrr}
\caption{Tabla completa de profesores para MAT1022, anio 2022, sesion PRIMAVERA.}\label{tab:prof-period-mat1022-2022-primavera-especializado}\\
\toprule
Profesor & Total & Obj. & Obj. \% & Bench. & Bench. \% & Cal. & DMU & GA-GB & Rank \\
\midrule
\endfirsthead
\toprule
Profesor & Total & Obj. & Obj. \% & Bench. & Bench. \% & Cal. & DMU & GA-GB & Rank \\
\midrule
\endhead
24085 & 42 & 42 & 100.0 & 0 & 0.0 & 9.84 & 63.17 & 73.36 & 1 \\
23425 & 40 & 40 & 100.0 & 1 & 2.5 & 9.36 & 63.83 & 70.78 & 2 \\
21853 & 20 & 20 & 100.0 & 0 & 0.0 & 9.29 & 72.67 & 75.42 & 3 \\
23452 & 20 & 20 & 100.0 & 0 & 0.0 & 9.52 & 79.67 & 74.58 & 4 \\
22521 & 14 & 14 & 100.0 & 0 & 0.0 & 8.74 & 70.00 & 73.66 & 5 \\
22446 & 44 & 43 & 97.7 & 1 & 2.3 & 9.59 & 59.55 & 70.83 & 6 \\
23897 & 40 & 39 & 97.5 & 1 & 2.5 & 9.59 & 72.67 & 71.20 & 7 \\
23411 & 60 & 57 & 95.0 & 0 & 0.0 & 8.97 & 62.56 & 72.12 & 8 \\
23453 & 33 & 31 & 93.9 & 1 & 3.0 & 8.24 & 57.98 & 71.02 & 9 \\
23836 & 12 & 11 & 91.7 & 1 & 8.3 & 8.88 & 66.67 & 66.15 & 10 \\
23148 & 11 & 10 & 90.9 & 0 & 0.0 & 8.05 & 53.94 & 65.15 & 11 \\
21192 & 43 & 39 & 90.7 & 0 & 0.0 & 9.07 & 68.99 & 66.91 & 12 \\
18269 & 36 & 32 & 88.9 & 1 & 2.8 & 8.62 & 71.11 & 75.64 & 13 \\
23835 & 20 & 17 & 85.0 & 0 & 0.0 & 7.61 & 78.00 & 78.75 & 14 \\
21779 & 9 & 6 & 66.7 & 0 & 0.0 & 6.67 & 73.33 & 71.76 &  \\
\bottomrule
\end{longtable}
\endgroup

\subsubsection{Anio 2023, sesion OTOÑO}

\begingroup
\scriptsize
\setlength{\tabcolsep}{3pt}
\begin{longtable}{lrrrrrrrrr}
\caption{Tabla completa de profesores para MAT1022, anio 2023, sesion OTOÑO.}\label{tab:prof-period-mat1022-2023-oto-o-especializado}\\
\toprule
Profesor & Total & Obj. & Obj. \% & Bench. & Bench. \% & Cal. & DMU & GA-GB & Rank \\
\midrule
\endfirsthead
\toprule
Profesor & Total & Obj. & Obj. \% & Bench. & Bench. \% & Cal. & DMU & GA-GB & Rank \\
\midrule
\endhead
23852 & 7 & 7 & 100.0 & 0 & 0.0 & 8.96 & 69.52 & 66.37 & 1 \\
23452 & 8 & 7 & 87.5 & 0 & 0.0 & 8.49 & 55.83 & 55.99 & 2 \\
21192 & 5 & 4 & 80.0 & 0 & 0.0 & 8.96 & 54.67 & 48.33 & 3 \\
22114 & 4 & 4 & 100.0 & 0 & 0.0 & 8.70 & 53.33 & 48.96 &  \\
23887 & 3 & 3 & 100.0 & 0 & 0.0 & 8.87 & 31.11 & 58.33 &  \\
24086 & 2 & 1 & 50.0 & 0 & 0.0 & 5.25 & 20.00 & 25.00 &  \\
23152 & 3 & 0 & 0.0 & 0 & 0.0 & 2.67 & 42.22 & 61.11 &  \\
22115 & 2 & 0 & 0.0 & 0 & 0.0 & 3.90 & 46.67 & 62.50 &  \\
23823 & 2 & 0 & 0.0 & 0 & 0.0 & 2.15 & 73.33 & 66.67 &  \\
\bottomrule
\end{longtable}
\endgroup

\subsubsection{Anio 2023, sesion PRIMAVERA}

\begingroup
\scriptsize
\setlength{\tabcolsep}{3pt}
\begin{longtable}{lrrrrrrrrr}
\caption{Tabla completa de profesores para MAT1022, anio 2023, sesion PRIMAVERA.}\label{tab:prof-period-mat1022-2023-primavera-especializado}\\
\toprule
Profesor & Total & Obj. & Obj. \% & Bench. & Bench. \% & Cal. & DMU & GA-GB & Rank \\
\midrule
\endfirsthead
\toprule
Profesor & Total & Obj. & Obj. \% & Bench. & Bench. \% & Cal. & DMU & GA-GB & Rank \\
\midrule
\endhead
23885 & 42 & 42 & 100.0 & 1 & 2.4 & 9.30 & 67.14 & 71.23 & 1 \\
23897 & 25 & 25 & 100.0 & 0 & 0.0 & 9.44 & 79.47 & 77.75 & 2 \\
23852 & 22 & 22 & 100.0 & 0 & 0.0 & 9.34 & 78.79 & 74.05 & 3 \\
20792 & 19 & 19 & 100.0 & 1 & 5.3 & 8.31 & 56.14 & 71.60 & 4 \\
23148 & 19 & 19 & 100.0 & 0 & 0.0 & 8.71 & 59.30 & 70.83 & 5 \\
22473 & 26 & 25 & 96.2 & 0 & 0.0 & 9.03 & 69.23 & 70.99 & 6 \\
23149 & 26 & 25 & 96.2 & 0 & 0.0 & 9.22 & 66.41 & 73.48 & 7 \\
24827 & 24 & 23 & 95.8 & 0 & 0.0 & 9.32 & 75.00 & 70.75 & 8 \\
24085 & 18 & 16 & 88.9 & 2 & 11.1 & 9.96 & 52.59 & 69.79 & 9 \\
22446 & 17 & 15 & 88.2 & 0 & 0.0 & 8.66 & 72.16 & 69.73 & 10 \\
22521 & 14 & 11 & 78.6 & 0 & 0.0 & 8.17 & 66.19 & 79.91 & 11 \\
23836 & 16 & 12 & 75.0 & 0 & 0.0 & 7.21 & 65.83 & 67.32 & 12 \\
21779 & 23 & 17 & 73.9 & 0 & 0.0 & 8.10 & 65.22 & 64.67 & 13 \\
23411 & 6 & 6 & 100.0 & 0 & 0.0 & 9.43 & 77.78 & 72.22 &  \\
\bottomrule
\end{longtable}
\endgroup

\subsubsection{Anio 2024, sesion OTOÑO}

\begingroup
\scriptsize
\setlength{\tabcolsep}{3pt}
\begin{longtable}{lrrrrrrrrr}
\caption{Tabla completa de profesores para MAT1022, anio 2024, sesion OTOÑO.}\label{tab:prof-period-mat1022-2024-oto-o-especializado}\\
\toprule
Profesor & Total & Obj. & Obj. \% & Bench. & Bench. \% & Cal. & DMU & GA-GB & Rank \\
\midrule
\endfirsthead
\toprule
Profesor & Total & Obj. & Obj. \% & Bench. & Bench. \% & Cal. & DMU & GA-GB & Rank \\
\midrule
\endhead
21192 & 10 & 10 & 100.0 & 2 & 20.0 & 9.37 & 58.00 & 51.67 & 1 \\
25503 & 2 & 2 & 100.0 & 0 & 0.0 & 8.50 & 50.00 & 58.33 &  \\
23452 & 1 & 1 & 100.0 & 0 & 0.0 & 9.10 & 66.67 & 41.67 &  \\
22473 & 5 & 4 & 80.0 & 0 & 0.0 & 7.48 & 65.33 & 53.33 &  \\
23148 & 5 & 4 & 80.0 & 0 & 0.0 & 7.90 & 60.00 & 49.58 &  \\
21854 & 6 & 4 & 66.7 & 0 & 0.0 & 6.43 & 55.56 & 52.43 &  \\
24086 & 5 & 2 & 40.0 & 0 & 0.0 & 6.38 & 53.33 & 54.17 &  \\
23152 & 3 & 0 & 0.0 & 0 & 0.0 & 3.07 & 80.00 & 59.72 &  \\
22114 & 1 & 0 & 0.0 & 0 & 0.0 & 1.10 & 53.33 & 50.00 &  \\
\bottomrule
\end{longtable}
\endgroup

\subsubsection{Anio 2024, sesion PRIMAVERA}

\begingroup
\scriptsize
\setlength{\tabcolsep}{3pt}
\begin{longtable}{lrrrrrrrrr}
\caption{Tabla completa de profesores para MAT1022, anio 2024, sesion PRIMAVERA.}\label{tab:prof-period-mat1022-2024-primavera-especializado}\\
\toprule
Profesor & Total & Obj. & Obj. \% & Bench. & Bench. \% & Cal. & DMU & GA-GB & Rank \\
\midrule
\endfirsthead
\toprule
Profesor & Total & Obj. & Obj. \% & Bench. & Bench. \% & Cal. & DMU & GA-GB & Rank \\
\midrule
\endhead
23452 & 29 & 29 & 100.0 & 0 & 0.0 & 9.39 & 79.31 & 72.05 & 1 \\
24097 & 29 & 29 & 100.0 & 0 & 0.0 & 9.53 & 80.92 & 60.70 & 2 \\
23885 & 28 & 28 & 100.0 & 0 & 0.0 & 9.15 & 66.43 & 70.31 & 3 \\
23823 & 17 & 17 & 100.0 & 0 & 0.0 & 9.31 & 76.47 & 60.78 & 4 \\
23978 & 15 & 15 & 100.0 & 0 & 0.0 & 9.20 & 68.00 & 65.69 & 5 \\
21779 & 34 & 33 & 97.1 & 1 & 2.9 & 8.89 & 70.00 & 71.45 & 6 \\
23836 & 55 & 53 & 96.4 & 2 & 3.6 & 9.55 & 63.88 & 62.65 & 7 \\
23852 & 55 & 53 & 96.4 & 0 & 0.0 & 9.19 & 68.36 & 65.87 & 8 \\
24086 & 26 & 25 & 96.2 & 0 & 0.0 & 8.89 & 71.28 & 74.68 & 9 \\
23152 & 22 & 21 & 95.5 & 0 & 0.0 & 9.00 & 74.24 & 71.40 & 10 \\
23153 & 21 & 20 & 95.2 & 0 & 0.0 & 8.90 & 64.44 & 64.98 & 11 \\
21262 & 19 & 18 & 94.7 & 0 & 0.0 & 8.25 & 74.39 & 69.08 & 12 \\
23148 & 19 & 18 & 94.7 & 0 & 0.0 & 8.77 & 71.93 & 69.08 & 13 \\
24085 & 46 & 43 & 93.5 & 1 & 2.2 & 9.98 & 63.62 & 65.31 & 14 \\
23411 & 27 & 25 & 92.6 & 0 & 0.0 & 9.03 & 74.32 & 66.28 & 15 \\
22446 & 13 & 12 & 92.3 & 0 & 0.0 & 8.28 & 64.62 & 69.39 & 16 \\
23839 & 22 & 20 & 90.9 & 0 & 0.0 & 8.37 & 58.18 & 65.72 & 17 \\
22471 & 14 & 12 & 85.7 & 0 & 0.0 & 8.00 & 63.33 & 67.26 & 18 \\
\bottomrule
\end{longtable}
\endgroup

\subsubsection{Anio 2025, sesion PRIMAVERA}

\begingroup
\scriptsize
\setlength{\tabcolsep}{3pt}
\begin{longtable}{lrrrrrrrrr}
\caption{Tabla completa de profesores para MAT1022, anio 2025, sesion PRIMAVERA.}\label{tab:prof-period-mat1022-2025-primavera-especializado}\\
\toprule
Profesor & Total & Obj. & Obj. \% & Bench. & Bench. \% & Cal. & DMU & GA-GB & Rank \\
\midrule
\endfirsthead
\toprule
Profesor & Total & Obj. & Obj. \% & Bench. & Bench. \% & Cal. & DMU & GA-GB & Rank \\
\midrule
\endhead
23836 & 39 & 39 & 100.0 & 0 & 0.0 & 9.58 & 78.12 & 74.41 & 1 \\
23885 & 26 & 26 & 100.0 & 0 & 0.0 & 9.35 & 76.67 & 73.40 & 2 \\
24085 & 26 & 26 & 100.0 & 0 & 0.0 & 9.55 & 63.85 & 65.38 & 3 \\
25503 & 20 & 20 & 100.0 & 0 & 0.0 & 9.80 & 69.00 & 68.33 & 4 \\
24097 & 16 & 16 & 100.0 & 0 & 0.0 & 9.68 & 88.75 & 74.87 & 5 \\
23153 & 15 & 15 & 100.0 & 0 & 0.0 & 9.69 & 78.67 & 70.42 & 6 \\
23411 & 10 & 10 & 100.0 & 0 & 0.0 & 9.57 & 81.33 & 74.17 & 7 \\
25531 & 10 & 10 & 100.0 & 0 & 0.0 & 9.87 & 70.67 & 71.67 & 8 \\
23852 & 41 & 40 & 97.6 & 1 & 2.4 & 9.06 & 75.28 & 71.24 & 9 \\
24086 & 35 & 34 & 97.1 & 0 & 0.0 & 9.37 & 72.95 & 76.19 & 10 \\
23978 & 34 & 33 & 97.1 & 0 & 0.0 & 9.42 & 74.71 & 66.91 & 11 \\
22446 & 19 & 18 & 94.7 & 0 & 0.0 & 9.28 & 81.05 & 58.99 & 12 \\
20792 & 17 & 16 & 94.1 & 0 & 0.0 & 9.52 & 85.10 & 79.41 & 13 \\
21779 & 31 & 28 & 90.3 & 0 & 0.0 & 8.53 & 71.40 & 72.31 & 14 \\
23452 & 20 & 18 & 90.0 & 0 & 0.0 & 9.46 & 76.00 & 63.54 & 15 \\
25843 & 23 & 20 & 87.0 & 0 & 0.0 & 8.01 & 71.30 & 63.59 & 16 \\
21262 & 8 & 8 & 100.0 & 0 & 0.0 & 8.21 & 68.33 & 64.58 &  \\
23152 & 8 & 8 & 100.0 & 0 & 0.0 & 8.49 & 62.50 & 69.79 &  \\
23823 & 3 & 3 & 100.0 & 0 & 0.0 & 9.77 & 64.44 & 55.56 &  \\
\bottomrule
\end{longtable}
\endgroup}{}}

\end{document}
